\section{Introduction}\label{sec:intro}

Adaptive systems rarely track a fixed target. In streaming estimation,
monitoring, and online decision-making, the data-generating distribution itself
drifts over time. A learner with finite memory must therefore answer a basic but
easy-to-miss question: how much of the past should still count once the world
that produced it has started to move away?

That question creates a structural trade-off. More past data reduces sampling
noise, but it also increases the age of the evidence used to describe the
present. Under drift, memory is simultaneously helpful and harmful: short
windows are fresh but noisy, while long windows are stable but stale. The paper
therefore studies a worst-case horizon problem, not a detector comparison.

For the finite-memory averaging class analyzed here, the tracking error splits
into a statistical term that decreases with memory and a staleness term that
increases with drift. In the linear staleness regime induced by the local
$W_2$-Lipschitz path bound, optimizing that balance yields a cube-root law:
$n^* \asymp (C_K/\zeta)^{2/3}$ and
$\mathcal{E}_{\min} \asymp C_K^{2/3}\zeta^{1/3}$. The same class also admits a
minimax lower bound in the critical-window regime, so the finite-memory floor is
structural rather than a quirk of one estimator.

The practical conflict is between two clocks. One clock measures the age at
which retained evidence stops being valid for the present. The other measures
when the accumulated data provide enough statistical evidence for a changepoint
alarm. Those clocks need not move together. The central conceptual distinction
is therefore simple: temporal validity is not changepoint evidence. A detector
can remain statistically silent while retained memory is already stale for the
present.

That gap motivates a conservative instrument rather than a stronger detector.
UMR is used here as a horizon regulator: it caps retained memory by a local
drift proxy so that the system does not wait for changepoint evidence before
discarding stale history. Its role in the paper is narrow and operational. It is
not presented as a universal detector and not as a claim that every backend must
adapt in the same way.

The contributions are fivefold:
\begin{itemize}
  \item It formalizes useful memory under drift as a variance-staleness
  trade-off in Wasserstein geometry.
  \item It proves a worst-case finite-memory horizon law under local
  $W_2$-Lipschitz drift.
  \item It shows a Gaussian-location lower bound establishing a structural
  tracking floor in the critical-window regime.
  \item It distinguishes temporal validity from changepoint evidence and
  identifies a cap-only regime.
  \item It instantiates the theory as a conservative horizon regulator and
  evaluates where it helps, saturates, and fails.
\end{itemize}

The empirical evidence follows that logic in stages. It first recovers the
predicted $U$-curve and the finite useful-memory optimum. It then shows the gap
between temporal validity and changepoint evidence, the cap-only regime, the
recovery asymmetry under alternating drift, and the nonlinear cost of horizon
misalignment. ELEC2 shows the same qualitative issue outside the synthetic
setting, while additional stability and sensitivity checks appear in the
appendix.

The theoretical results are stated in $W_2$; the Sinkhorn divergence is the
computational measurement layer used in the synthetic experiments. The window
size result of \citet{hinder2023window} motivates the broader issue, but the
contribution here is more specific: memory has a finite temporal validity horizon
under worst-case drift. The next step is correspondingly limited: characterize
narrower path classes that change the staleness law without diluting the current
worst-case statement.
